\documentclass[a4paper,11pt]{article}

\usepackage[utf8]{inputenc}
\usepackage[T1]{fontenc}
\usepackage[french]{babel}
\usepackage{amsmath,amssymb}
\usepackage{hyperref}
\usepackage{geometry}
\geometry{margin=2.5cm}

\title{Rapport sur le projet PySpark avec le jeu de données MovieLens}
\author{Moubarak LASSISSI}
\date{\today}

\begin{document}

\maketitle

\tableofcontents

\section{Introduction}

L'objectif de ce projet est de découvrir et de pratiquer PySpark,
l'API Python d'Apache Spark, à travers l'analyse du jeu de données
\emph{MovieLens}. Il couvre plusieurs aspects essentiels
du travail avec Spark~: exploration de données, agrégations,
SQL distribué, jointures, traitement de chaînes (genres), apprentissage
automatique  et notions de performance (partitions,
cache, plans d'exécution).

\section{Présentation de PySpark}
Apache Spark est un moteur de calcul distribué pour le Big Data.
Il permet de traiter de grands volumes de données en mémoire, sur
un cluster, en fournissant un modèle de programmation simple et
expressif tels que :
\begin{itemize}
  \item le calcul distribué sur plusieurs machines (ou cœurs),
  \item la tolérance aux pannes grâce au modèle de résilience,
  \item la possibilité de chaîner des opérations complexes
        (filtres, agrégations, jointures, ML, SQL, etc.).
\end{itemize}

\subsection{PySpark et les base des données}
PySpark est l'interface de Spark pour le langage Python, il est éssentiellenlement utilisé pour~:
\begin{itemize}
  \item manipuler des DataFrames distribués,
  \item exécuter des requêtes SQL sur de grands ensembles de données,
  \item appliquer des algorithmes d'apprentissage automatique
        via MLlib.
\end{itemize}

\section{Application}
\subsection{Jeu de données MovieLens}

Le jeu de données \emph{MovieLens} contient des évaluations de films
par des utilisateurs. Dans ce projet, nous avons utilisé la version
\emph{ml-latest-small}, composée notamment de deux fichiers CSV~:
\begin{itemize}
  \item \texttt{ratings.csv}~: contient les colonnes
        \texttt{userId}, \texttt{movieId}, \texttt{rating},
        \texttt{timestamp}.
  \item \texttt{movies.csv}~: contient les colonnes
        \texttt{movieId}, \texttt{title}, \texttt{genres}.
\end{itemize}

Ces fichiers ont été lus avec \texttt{spark.read.csv} en demandant
l'inférence de schéma et la présence d'un en-tête (option
\texttt{header=True, inferSchema=True}).

\section{Tâches 0 et 1~: mise en place et exploration}

\subsection{Tâche 0~: création de la SparkSession}

La première étape du notebook consiste à installer PySpark
(pour Colab) puis à créer une \texttt{SparkSession}~:
\begin{itemize}
  \item définition d'un nom d'application
        (\texttt{PySpark\_Student\_Exercises\_Movies}),
  \item mode \texttt{local[*]} pour utiliser tous les cœurs
        disponibles sur la machine locale.
\end{itemize}

Cette session permet ensuite de créer des DataFrames à partir
des fichiers CSV et d'exécuter des requêtes SQL.

\subsection{Tâche 1.1~: inspection des données}

Sur les DataFrames \texttt{ratings} et \texttt{movies},
les opérations suivantes ont été réalisées~:
\begin{itemize}
  \item affichage des 10 premières lignes avec \texttt{show(10)},
  \item affichage des schémas avec \texttt{printSchema()},
  \item comptage du nombre de lignes avec \texttt{count()}.
\end{itemize}

Ces réquêtes permettent de vérifier la structure des données, les types de
colonnes et la taille du jeu de données.

\subsection{Tâche 1.2~: sélection et transformation de colonnes}

Pour \texttt{ratings}, un nouveau DataFrame \texttt{ratings\_selected}
a été créé en~:
\begin{itemize}
  \item sélectionnant les colonnes \texttt{userId}, \texttt{movieId},
        \texttt{rating},
  \item renommant \texttt{movieId} en \texttt{film\_id},
  \item castant explicitement \texttt{rating} en \texttt{float}.
\end{itemize}

Ces opérations utilisent des fonctions PySpark telles que
\texttt{select}, \texttt{alias} et \texttt{cast}.

\subsection{Tâche 1.3~: filtrage}

Plusieurs exemples de filtrage ont été implémentés~:
\begin{itemize}
  \item toutes les évaluations de l'utilisateur \texttt{userId = 1},
  \item toutes les évaluations strictement supérieures à 4{,}5,
  \item toutes les évaluations dont \texttt{movieId} est dans l'ensemble
        $\{1, 50, 100\}$ en utilisant la méthode \texttt{isin()}.
\end{itemize}

Cela illustre l'utilisation de \texttt{filter} (ou \texttt{where})
et des expressions de colonnes.

\subsection{Tâche 1.4~: tri et limitation}

Toujours sur \texttt{ratings}, les opérations suivantes ont été réalisées~:
\begin{itemize}
  \item affichage des 20 plus hautes évaluations (tri par
        \texttt{rating} décroissant),
  \item affichage des 10 plus basses évaluations de l'utilisateur
        \texttt{600},
  \item tri combiné par \texttt{rating} décroissant puis
        \texttt{timestamp} croissant.
\end{itemize}

On utilise \texttt{orderBy} avec des colonnes ordonnées
ascendantes ou descendantes.

\subsection{Tâche 1.5~: colonnes dérivées}

De nouvelles colonnes ont été créées~:
\begin{itemize}
  \item \texttt{rating\_x2 = rating * 2},
  \item \texttt{positive\_rating}~: vaut 1 si \texttt{rating} $\geq 4$,
        sinon 0 (via \texttt{when} et \texttt{otherwise}),
  \item \texttt{log\_rating}~: $\log_{10}(rating + 1)$.
\end{itemize}

Cette étape montre comment enrichir un DataFrame avec des
colonnes calculées.

\subsection{Tâche 1.6~: valeurs manquantes}

Même si le jeu de données MovieLens ne contient pas de valeurs
manquantes, une colonne artificielle \texttt{fake\_null\_col} a
été ajoutée avec des valeurs nulles, puis~:
\begin{itemize}
  \item remplissage des valeurs nulles avec \texttt{fillna},
  \item remplacement conditionnel de certaines valeurs de \texttt{rating}
        (par exemple transformer les notes 5 en 10) avec \texttt{when},
  \item suppression de lignes contenant des valeurs nulles avec
        \texttt{dropna}.
\end{itemize}

\subsection{Tâche 1.7~: distinct et déduplication}

Des opérations de déduplication ont été menées~:
\begin{itemize}
  \item comptage du nombre d'utilisateurs distincts,
  \item comptage du nombre de films distincts notés,
  \item recherche d'éventuels doublons sur les paires
        \texttt{(userId, movieId)},
  \item suppression potentielle de doublons avec
        \texttt{dropDuplicates}.
\end{itemize}

\section{Tâche 2~: agrégations et fonctions de fenêtre}

\subsection{Tâche 2.1~: agrégations simples}

Sur le DataFrame \texttt{ratings}, les agrégations suivantes
ont été calculées~:
\begin{itemize}
  \item minimum, maximum et moyenne des notes,
  \item nombre total d'évaluations,
  \item nombre d'évaluations par \texttt{movieId}.
\end{itemize}

Cela utilise les fonctions d'agrégation \texttt{min}, \texttt{max},
\texttt{avg}, \texttt{count}.

\subsection{Tâche 2.2~: groupBy}

Les agrégations suivantes ont été calculées avec \texttt{groupBy}~:
\begin{itemize}
  \item moyenne et nombre de notes par film,
  \item moyenne et nombre de notes par utilisateur.
\end{itemize}

Ces résultats donnent une première idée de la popularité des films
et de l'activité des utilisateurs.

\subsection{Tâche 2.3~: top films}

Deux types de ``top'' ont été calculés~:
\begin{itemize}
  \item les 20 films les plus notés (en termes de nombre total
        d'évaluations),
  \item les 20 meilleurs films en moyenne, en imposant un seuil
        minimal de 50 évaluations pour éviter les biais dus à un
        faible nombre de notes.
\end{itemize}

Une jointure avec \texttt{movies} permet d'ajouter les titres des films.

\subsection{Tâche 2.4~: fonctions de fenêtre}

Des fonctions de fenêtre (\emph{window functions}) ont été utilisées
avec des partitions par \texttt{movieId} ou \texttt{userId}~:
\begin{itemize}
  \item classement des évaluations dans l'ordre chronologique
        (\texttt{row\_number} par film),
  \item colonne de décalage (\texttt{lag}) pour obtenir la note
        précédente d'un même utilisateur,
  \item moyenne de la note calculée en fenêtre par film
        (moyenne glissante).
\end{itemize}

Ces fonctionnalités sont utiles pour analyser des données séquentielles.

\section{Tâche 3~: Spark SQL}

\subsection{Vues temporaires}

Les DataFrames \texttt{ratings} et \texttt{movies} ont été
enregistrés comme vues temporaires \texttt{ratings\_table}
et \texttt{movies\_table} grâce à
\texttt{createOrReplaceTempView}. Cela permet d'écrire des
requêtes SQL directement sur ces tables.

\subsection{Tâche 3.1~: requêtes SQL simples}

Plusieurs requêtes SQL simples ont été exécutées~:
\begin{itemize}
  \item affichage des 20 premières lignes,
  \item comptage du nombre total de notes,
  \item comptage d'utilisateurs distincts,
  \item calcul de la note moyenne globale.
\end{itemize}

\subsection{Tâche 3.2~: regroupements et classement}

En SQL, les mêmes agrégations que précédemment ont été
reproduites~:
\begin{itemize}
  \item moyenne et nombre de notes par film (avec jointure
        pour récupérer le titre),
  \item meilleurs et pires films sous contrainte d'un nombre
        minimal de notes (par exemple au moins 100),
  \item utilisateurs les plus actifs (ceux qui notent le plus).
\end{itemize}

\subsection{Tâche 3.3~: CASE WHEN et catégories de notes}

Une catégorisation des notes a été implémentée en SQL avec
\texttt{CASE WHEN}~:
\begin{itemize}
  \item notes $\geq 4.0$~: catégorie \texttt{high},
  \item notes entre 3.0 et 3.9~: catégorie \texttt{medium},
  \item notes $< 3.0$~: catégorie \texttt{low}.
\end{itemize}

Une distribution du nombre de notes par catégorie a ensuite été
calculée.

\subsection{Tâche 3.4~: fenêtres en SQL}

Des fonctions de fenêtre SQL ont été utilisées pour~:
\begin{itemize}
  \item obtenir, pour chaque genre, le top 10 des films selon leur
        note moyenne (avec explosion de la colonne \texttt{genres}),
  \item déterminer la première note de chaque utilisateur
        (via \texttt{ROW\_NUMBER} et filtrage),
  \item calculer une moyenne glissante de la note par film en fonction
        du temps.
\end{itemize}

\section{Tâche 4~: jointures et analyse des genres}

\subsection{Tâche 4.1~: jointure de base}

Un DataFrame \texttt{ratings\_movies} a été créé en joignant
\texttt{ratings} et \texttt{movies} sur \texttt{movieId}. On a
ainsi pu afficher des triplets \texttt{(userId, title, rating)}
et effectuer des analyses combinant notes et informations sur
les films.

\subsection{Tâche 4.2~: parsing et explosion des genres}

La colonne \texttt{genres} est une chaîne de caractères de type
\texttt{Action|Comedy|Drama}. Les étapes suivantes ont été réalisées~:
\begin{itemize}
  \item séparation de la chaîne en tableau avec \texttt{split},
  \item explosion du tableau en plusieurs lignes avec \texttt{explode},
  \item calcul du nombre de notes et de la note moyenne par genre.
\end{itemize}

\subsection{Tâche 4.3~: agrégations par genre, titre et année}

À partir du DataFrame explosé sur les genres~:
\begin{itemize}
  \item calcul de la note moyenne par genre,
  \item calcul du nombre de notes et de la note moyenne par film
        (titre),
  \item extraction de l'année depuis le titre au format
        \texttt{Titre (YYYY)} à l'aide d'une expression régulière,
  \item calcul du nombre de notes et de la note moyenne par
        couple (genre, année).
\end{itemize}

\subsection{Tâche 4.4~: left anti join}

Une jointure \emph{left anti} entre \texttt{movies} et \texttt{ratings}
sur \texttt{movieId} a permis d'identifier les films présents dans
\texttt{movies.csv} mais n'ayant reçu aucune note. Le nombre de ces
films a été compté.

\section{Tâche 5~: classification binaire avec MLlib}

\subsection{Tâche 5.1~: création du label}

L'objectif est de prédire si une note sera $\geq 4.0$ ou non.
Pour cela, une colonne \texttt{label} a été ajoutée au DataFrame
\texttt{ratings}~:
\[
  \text{label} =
  \begin{cases}
    1 & \text{si } \text{rating} \geq 4.0, \\
    0 & \text{sinon.}
  \end{cases}
\]

\subsection{Tâche 5.2~: features et VectorAssembler}

Plusieurs variables ont été utilisées comme \emph{features}~:
\begin{itemize}
  \item la note \texttt{rating} ,
  \item le \texttt{timestamp} brut,
  \item un \texttt{timestamp} normalisé entre 0 et 1,
  \item le nombre total de notes données par chaque utilisateur
        (\texttt{user\_rating\_count}),
  \item un \texttt{log\_timestamp} comme transformation logarithmique
        du temps.
\end{itemize}

La normalisation du \texttt{timestamp} se fait en utilisant le
minimum et le maximum des timestamps dans le jeu de données,
puis en calculant~:
\[
  \text{norm\_timestamp} =
  \frac{\text{timestamp} - \text{min\_ts}}{\text{max\_ts} - \text{min\_ts}}.
\]

Ensuite, un \texttt{VectorAssembler} regroupe ces colonnes
numériques en un unique vecteur de features, stocké dans une
colonne \texttt{features}.

\subsection{Tâche 5.3~: séparation train / test}

Le DataFrame des features et du label a été séparé aléatoirement
en deux ensembles~:
\begin{itemize}
  \item 70\% pour l'entraînement (train),
  \item 30\% pour le test.
\end{itemize}

\subsection{Tâche 5.4~: entraînement d'une régression logistique}

Un modèle de \texttt{LogisticRegression} a été entraîné sur
l'ensemble d'entraînement. Après l'entraînement~:
\begin{itemize}
  \item les coefficients et l'interception (biais) du modèle
        ont été affichés,
  \item le modèle a été appliqué à l'ensemble de test pour
        produire des prédictions et des probabilités.
\end{itemize}

\subsection{Tâche 5.5~: évaluation du modèle}

Les performances de la régression logistique ont été évaluées
avec plusieurs métriques~:
\begin{itemize}
  \item l'aire sous la courbe ROC (AUC) via
        \texttt{BinaryClassificationEvaluator},
  \item la matrice de confusion (TP, FP, TN, FN) calculée à partir
        des prédictions,
  \item l'accuracy, la précision et le rappel en fonction de ces
        quantités.
\end{itemize}

Ces métriques donnent une vision globale de la qualité de la
classification.

\subsection{Tâche 5.6~: amélioration du modèle}

Pour améliorer et comparer les performances, d'autres modèles
ont été entraînés~:
\begin{itemize}
  \item un \texttt{DecisionTreeClassifier} (arbre de décision),
  \item un \texttt{RandomForestClassifier} (forêt aléatoire),
\end{itemize}

en utilisant les mêmes features. Les AUC correspondants ont été
calculés, permettant de comparer les différents algorithmes.

\section{Tâche 6~: performance et exécution}

\subsection{Tâche 6.1~: partitions}

Le nombre de partitions des DataFrames \texttt{ratings} et
\texttt{movies} a été inspecté via
\texttt{rdd.getNumPartitions()}. Ce nombre contrôle le degré
de parallélisme des traitements.

\subsection{Tâche 6.2~: repartition et coalesce}

Deux opérations ont été comparées~:
\begin{itemize}
  \item \texttt{repartition(n)}~: redistribue les données, implique
        généralement un \emph{shuffle} ,
  \item \texttt{coalesce(n)}~: réduit le nombre de partitions avec
        peu ou pas de \emph{shuffle}.
\end{itemize}

Les plans physiques de ces opérations ont été observés avec
\texttt{explain(True)} pour voir l'impact sur les étapes
d'exécution.

\subsection{Tâche 6.3~: cache et temps d'exécution}

Le DataFrame \texttt{ratings} a été mis en cache (et
repartitionné)~:
\begin{itemize}
  \item mesure du temps de \texttt{count()} sur la première exécution,
  \item mesure du temps sur une seconde exécution (qui profite du cache),
  \item même type de mesure pour le calcul de la moyenne de note
        par film, calculée deux fois.
\end{itemize}

Cela montre l'intérêt du cache pour des calculs répétés sur les
mêmes données.

\subsection{Tâche 6.4~: analyse des plans d'exécution}

La méthode \texttt{explain(True)} a été utilisée pour analyser
les plans physiques~:
\begin{itemize}
  \item d'une jointure entre \texttt{ratings} et \texttt{movies},
  \item d'un \texttt{groupBy} avec agrégation,
  \item d'une fonction de fenêtre (moyenne de notes par film).
\end{itemize}

Ces plans mettent en évidence les frontières de \emph{shuffle}
et les différentes étapes (lecture, projection, agrégation,
jointure), ce qui est essentiel pour comprendre et
optimiser les performances de Spark.

\section{Conclusion}

Ce projet a permis de mettre en pratique un large éventail de
fonctionnalités de PySpark sur un jeu de données réel~:
\begin{itemize}
  \item chargement et exploration de données massives avec les
        DataFrames,
  \item filtrage, tri, colonnes dérivées et gestion des valeurs
        manquantes,
  \item agrégations avancées et fonctions de fenêtre,
  \item utilisation de Spark SQL pour effectuer des requêtes
        déclaratives,
  \item jointures et traitement de chaînes (genres de films),
  \item mise en place d'un problème de classification binaire avec
        MLlib et comparaison de plusieurs modèles,
  \item introduction aux notions de partitions, cache et analyse
        des plans d'exécution.
\end{itemize}

\end{document}
